\documentclass[11pt]{article}

\usepackage[margin=1in]{geometry}
\usepackage{amsmath, amssymb}
\usepackage{booktabs}
\usepackage{hyperref}
\usepackage{enumitem}
\usepackage{listings}
\usepackage{xcolor}
\usepackage{framed}
\usepackage{caption}
\usepackage{array}
% Pull labels from the main paper so cross-references resolve.
% Compile ../paper/stateval.tex first; xr reads its .aux file to recover labels.
\usepackage{xr}
\externaldocument[main-]{../paper/stateval}

\hypersetup{
    colorlinks=true,
    linkcolor=blue!50!black,
    citecolor=blue!50!black,
    urlcolor=blue!50!black
}

\definecolor{shadecolor}{rgb}{0.96,0.96,0.92}

\title{\textbf{Companion Workbook to ``How to Do Statistical Evaluations in
ECE/CS Papers''}\\
\large Project-Type Tables, an Evaluation-Plan Worksheet, an Autopsy of a
Fragile Evaluation, and Exercises}

\author{Bhaskar Krishnamachari\\
University of Southern California\\
\texttt{bkrishna@usc.edu}}

\date{}

\begin{document}

\maketitle

\begin{abstract}
This is a companion workbook for the tutorial \emph{How to Do Statistical
Evaluations in ECE/CS Papers} (the ``main paper''). It contains the
student-facing scaffolding that did not fit in the main paper: a project-type
translation table for matching common ECE/CS evaluation patterns to
appropriate methods, an evaluation-plan worksheet to fill out before running
experiments, an autopsy of a typical fragile evaluation and how to fix it,
and a set of exercises that map to the main paper's sections. It is intended
to be printed out, marked up, and used as a teaching artifact for student
research projects. References to sections beginning with \S\ldots refer to
the main paper.
\end{abstract}

\section{Choose Your Evaluation Type}
\label{sec:project-type}

ECE/CS students work on a wide range of empirical projects, and the unit of
analysis, the dominant metric, the most common mistake, and the recommended
method differ noticeably across them. Table~\ref{tab:project-type} lays out
the most common project types side-by-side. Students should locate
themselves in the table before reading the main paper, and reread the row
that applies to their project before each milestone.

\begin{table}[h]
\centering
\footnotesize
\setlength{\tabcolsep}{4pt}
\renewcommand{\arraystretch}{1.25}
\begin{tabular}{p{0.18\linewidth}p{0.16\linewidth}p{0.18\linewidth}p{0.20\linewidth}p{0.20\linewidth}}
\toprule
Project type & Unit of analysis & Typical metrics & Common mistake & Recommended default \\
\midrule
Wireless / network simulation & seed / topology / workload & delay, PDR, throughput, energy & treating packets as independent & paired seeds, bootstrap CIs, multiple regimes \\
Systems benchmark & run / machine / configuration & latency, throughput, $p_{95}$, $p_{99}$ & ignoring warmup, order, thermal effects & randomized order, warmup discarded, CDFs, CIs \\
ML classifier & held-out instance / group / time window & accuracy, AP, $F_1$, AUROC, calibration & tuning the threshold on the test set & validation/test split, AP for imbalance, paired bootstrap \\
Scheduler / optimizer & workload instance / seed & makespan, latency, utilization & only testing easy workloads & factorial regimes, ablations, paired tests \\
Embedded / IoT experiment & device / run / environment & energy, packet loss, deadline misses & hidden hardware or environment confounders & blocking, calibration, repeated trials \\
Simulation paper & scenario / seed & depends on the model & trusting simulator output blindly & sanity checks, conservation, baseline reproduction \\
\bottomrule
\end{tabular}
\caption{Project-type translation table. Find the row that matches your
project and use it as the lens through which to read the main paper. The
``unit of analysis'' column is what the main paper calls the unit of
analysis (\S4.1 of the main paper); the ``recommended default'' column is
the minimum-viable method, not the only acceptable one.}
\label{tab:project-type}
\end{table}

\paragraph{Translation boxes.} Three short translations of the main paper's
running scheduling example to other ECE/CS subfields.

\begin{shaded*}
\noindent\textbf{Wireless networking.}
Unit of analysis: topology / seed.
Primary metric: packet delivery ratio (PDR), latency, routing overhead, energy.
Pitfall: treating packets as independent (\S4.1).
Good plot: PDR vs.\ mobility speed with paired bootstrap CIs.
\end{shaded*}

\begin{shaded*}
\noindent\textbf{ML / anomaly detection.}
Unit of analysis: held-out flow / device / user / time window.
Primary metric: AP (under class imbalance), AUROC, recall at fixed
false-positive rate.
Pitfall: leakage across time or devices.
Good plot: precision--recall curve with calibration plot.
\end{shaded*}

\begin{shaded*}
\noindent\textbf{Embedded systems.}
Unit of analysis: repeated run / device / environment block.
Primary metric: energy, latency, missed deadlines.
Pitfall: temperature, battery state, background tasks (main \S\ref{main-sec:randomization}).
Good plot: energy vs.\ latency tradeoff.
\end{shaded*}

\section{Anatomy of a Fragile Evaluation}
\label{sec:autopsy}

A fragile evaluation often \emph{looks} fine on first read. The headline
sounds confident, the numbers are big, and the statistical machinery is
either absent or perfunctory. Below is a typical example, lightly
fictionalized, of the kind of single-line claim that recurs in beginning
student papers. We then take it apart and rewrite it.

\begin{shaded*}
\noindent\emph{``Our protocol improves throughput by 35\%.''}
\end{shaded*}

The sentence has at least nine separate problems, none of which are
visible on the surface:

\begin{enumerate}[leftmargin=*, itemsep=1pt]
    \item One topology was used.
    \item One seed was used.
    \item The baseline is a strawman version of the closest published method.
    \item No confidence interval is reported, on the difference or on either arm.
    \item Only throughput is reported; delay, energy, fairness, and overhead are silent.
    \item Parameters were tuned on the proposed method but not on the baseline.
    \item Mobility, failure, or workload stress conditions were not tested.
    \item The simulator was not validated against any analytical or published baseline.
    \item Per-packet measurements were treated as independent; the effective $n$ is one.
\end{enumerate}

A defensible rewrite of the same claim, using the main paper's machinery,
might read:

\begin{shaded*}
\noindent\emph{``Across 30 random seeds and 5 topologies (factorial design,
$5\times 30 = 150$ paired runs), our protocol increases mean throughput by
35\% (95\% bootstrap CI: 31--39\%) over the strongest published baseline at
moderate offered load $\rho = 0.6$. The improvement holds in 4 of 5
topologies; in the fifth (a sparse line topology) the protocols are within
1\%. Routing overhead increases by 12\% (CI: 9--15\%); energy is within 2\%.
Tuning was performed on a held-out validation topology; the test
topologies above were not used during tuning. Simulator delay and
queue-length outputs match the analytical M/M/1 model within 1.5\% in the
no-contention limit.''}
\end{shaded*}

The rewrite is longer, but each clause closes a specific reviewer
objection. The original sentence closed none.

\section{Evaluation Plan Worksheet}
\label{sec:worksheet}

Fill out this worksheet \emph{before} running experiments. The intent is
to catch ambiguity in the design phase, when changing the plan is cheap.
The last item, in particular, separates a real hypothesis from a wish.

\newcommand{\wsline}{\vspace{6pt}\hrule height 0.4pt\vspace{6pt}\hrule height 0.4pt\vspace{8pt}}

\begin{enumerate}[leftmargin=*, itemsep=2pt]
    \item Main claim (one sentence, plain English):\wsline
    \item Falsifiable hypothesis:\wsline
    \item Primary metric:\wsline
    \item Secondary metrics (including cost/overhead):\wsline
    \item Unit of analysis (seed / trace / topology / user / \ldots):\wsline
    \item Baselines (and why they are the strongest reasonable ones):\wsline
    \item Independent variables / factors:\wsline
    \item Regimes to test (load, distribution, scale, mobility, \ldots):\wsline
    \item Number of seeds / runs and a one-line justification:\wsline
    \item Pairing / blocking plan (CRN? same workloads? blocked by topology?):\wsline
    \item Statistical summary (mean / median / $p_{99}$ / AP / \ldots):\wsline
    \item Statistical test, if any (paired Wilcoxon / Friedman / McNemar / bootstrap / none):\wsline
    \item Figures to produce (CDF / interaction / Pareto / paired-difference / \ldots):\wsline
    \item Ablations:\wsline
    \item Stress tests, regimes where the method may fail:\wsline
    \item Simulator / measurement validation steps:\wsline
    \item Reproducibility artifacts (code, seeds, traces, hardware):\wsline
    \item \emph{What result would make us weaken or abandon the claim?}\wsline
\end{enumerate}

The last item is the one that distinguishes a real empirical hypothesis
from a wish, and it is the one students most often skip. If no
experimental outcome would change the claim, the claim is not falsifiable
and the experiment is decoration.

\subsection*{A filled-out example: FIFO vs.\ SRPT}

For students who have not seen what level of specificity to aim for, the
following is the worksheet completed for the running example used in the
main paper.

\begin{enumerate}[leftmargin=*, itemsep=1pt]
    \item Main claim: SRPT lowers mean job completion latency relative to FIFO under heavy-tailed workloads, with the gap growing in offered load.
    \item Falsifiable hypothesis: at $\rho \geq 0.5$ under truncated-Pareto job sizes, mean per-seed latency under SRPT is lower than under FIFO.
    \item Primary metric: per-seed mean job completion latency.
    \item Secondary metrics: per-seed $p_{99}$ latency; per-seed throughput; preemption count (overhead).
    \item Unit of analysis: \emph{seed}. Each seed redraws arrival times and job sizes; per-seed values are the independent observations for CIs and tests.
    \item Baselines: FIFO (a reasonable floor) and a published shortest-job-first variant (the strongest reasonable point of comparison; not a strawman).
    \item Independent variables: scheduling policy (FIFO, SRPT); offered load $\rho$; job-size distribution (truncated Pareto with $\alpha \in \{1.5, 2.5\}$, exponential).
    \item Regimes: $\rho \in \{0.3, 0.5, 0.7, 0.9\}$ crossed with the three job-size distributions; full $4 \times 3$ grid.
    \item Number of seeds: 30 per cell; chosen to give 80\% power against a Cohen's $d = 0.5$ at $\alpha = 0.05$ (paired).
    \item Pairing / blocking plan: common random numbers (one workload per seed, replayed under each policy); per-seed differences are the analysis unit.
    \item Statistical summary: per-cell mean of paired differences (FIFO $-$ SRPT) with 95\% bootstrap CI on the difference; $p_{99}$ reported per cell.
    \item Statistical test: paired Wilcoxon signed-rank on per-seed latency differences; significance reported alongside, not in place of, the CI.
    \item Figures: latency CCDF on log--log axes; interaction plot of mean latency vs.\ $\rho$ with paired bootstrap CIs; bootstrap distribution of the per-seed difference.
    \item Ablations: replace SRPT with shortest-job-first (no preemption) to isolate the contribution of preemption.
    \item Stress tests: very high load ($\rho = 0.95$); exponential job sizes (where the SRPT advantage should shrink); very small job buffers.
    \item Validation: simulator passes M/M/1 baseline reproduction in the exponential-size limit; conservation laws (jobs in $=$ jobs completed $+$ jobs in flight) hold to machine precision.
    \item Reproducibility: \texttt{figures.py} regenerates every figure from a fixed seed range; full code, seeds, and dependency versions in the project repository.
    \item What would weaken the claim? If the per-cell paired-difference CI overlaps zero at any $\rho \geq 0.5$ under truncated Pareto, or if the SRPT advantage at $\rho = 0.9$ is below 50\%, we would weaken the strength of the claim and report it.
\end{enumerate}

\section{Exercises}
\label{sec:exercises}

These exercises map roughly to the structure of the main paper. They are
short and meant to be done with paper, a calculator, and (where indicated)
a Python REPL. Solutions are not provided; they are the students' work.

\paragraph{E1. Strengthen a weak claim (main \S\ref{main-sec:weak-vs-strong}).} Rewrite
the following sentence as a strong evaluation claim, using the main
paper's reporting order:
\begin{quote}
\emph{``Our cache improves throughput by 25\%.''}
\end{quote}

\paragraph{E2. Identify the unit of analysis (main \S\ref{main-sec:unit-of-analysis}).} For each scenario
below, name the unit of analysis and the size of $n$.
\begin{enumerate}[label=(\alph*), leftmargin=*]
    \item One simulation run on one topology produces 10{,}000 packet
    delays.
    \item Twenty seeded simulation runs on one topology each produce
    10{,}000 packet delays per run.
    \item Ten topologies, with five seeded runs each, with 10{,}000 packet
    delays per run.
    \item A learned anomaly detector trained on traffic from 100 devices
    and evaluated on traffic from 20 held-out devices, with thousands of
    flows per device.
\end{enumerate}

\paragraph{E3. Plot choice (main \S\ref{main-sec:reporting}).} Pick the most informative plot for each:
(a) latency comparison of two schedulers across orders of magnitude;
(b) throughput vs.\ load for two protocols, with uncertainty;
(c) classifier comparison on imbalanced data;
(d) energy vs.\ latency for several configurations of a method.

\paragraph{E4. Paired or unpaired? (main \S\ref{main-sec:hypothesis-testing}).}
\begin{enumerate}[label=(\alph*), leftmargin=*]
    \item Two schedulers run on the same 30 seeded workloads.
    \item Two protocols run on different randomly drawn topologies.
    \item Two ML models evaluated on the same test set.
    \item Two algorithms evaluated on traces from disjoint pools of users.
\end{enumerate}
For each, state which test you would use by default and why.

\paragraph{E5. Compute an effect size and a CI (main \S\ref{main-sec:ci}, Listing~\ref{main-lst:bootstrap}).} You
are given 20 paired per-seed measurements of FIFO and SRPT mean latency.
Write the four-line numpy snippet that produces (i) the mean per-seed
reduction, (ii) the relative reduction in percent, and (iii) a 95\%
percentile bootstrap CI on the mean reduction.

\paragraph{E6. Find the flaw.}
\begin{quote}
\emph{``We ran one seed and report $p < 0.05$ over 10{,}000 packets.''}
\end{quote}
Identify every flaw in this sentence and propose a corrected version.

\paragraph{E7. Design a factorial experiment with a budget (main \S\ref{main-sec:multifactor}).} Assume
that one ``run'' means one policy executed on one workload, and that a
paired seed therefore costs two runs (one per policy). You have a
budget of 90 runs. Design a factorial experiment over (offered load $\in
\{0.3, 0.6, 0.9\}$) $\times$ (job-size distribution $\in \{\text{exponential}, \text{truncated Pareto}\}$)
for two policies, with paired seeds. How many complete paired seeds per
regime can you afford, and how many runs are left over? Which
interactions can you estimate from the resulting design?

\paragraph{E8. Reproducibility checklist (main \S\ref{main-sec:reproducibility}).} For a project of your choosing
in your group, write a short reproducibility statement that lists exactly
the information another student would need to re-run your experiments and
obtain the same numbers. Compare it against the main paper's
reproducibility section and Appendix~\ref{main-sec:checklist} (the
pre-submission checklist).

\paragraph{E9. Equivalence (main \S\ref{main-sec:equivalence}).} A new lightweight model claims to
match a heavy baseline. You measure mean test accuracy as 84.1\%
(baseline) vs.\ 83.7\% (proposed), with 95\% CI on the difference of
$[-1.0\%, +0.2\%]$. The team agrees on an equivalence margin of 1.5
percentage points. Does the experiment support the equivalence claim?
What if the margin were 0.5 percentage points instead?

\paragraph{E10. The autopsy, in your own area.} Find a published claim
in your subfield (or invent a plausible one) that has the same surface
shape as the autopsy in \S\ref{sec:autopsy} of this workbook. List the
problems and write a defensible rewrite.

\end{document}
